% ============================================================================
%  Capitolo 6 -- Rappresentazione dei numeri: i reali            (Lezione 06)
% ============================================================================
\chapter{Rappresentazione dei numeri: i reali}
\label{ch:reals}

\epigraphline{Conterò poco, è vero: --- diceva l'Uno ar Zero --- ma tu che
vali? Gnente: propio gnente.}{Trilussa, \emph{Nummeri}, 1944}

\begin{chapgoals}
Alla fine di questo capitolo saprete
\begin{itemize}
\item codificare interi con segno in complemento a uno e in eccesso-$k$, e dire a che cosa serve ciascuna codifica;
\item rappresentare una grandezza frazionaria in virgola fissa e calcolare con essa usando istruzioni intere;
\item spiegare mantissa, esponente e normalizzazione nella virgola mobile;
\item codificare e decodificare un numero in IEEE~754 a precisione singola;
\item riconoscere zero, numeri denormalizzati, infinito e NaN;
\item dire perché l'ESP32-C6 non ha istruzioni in virgola mobile e che cosa fanno invece i giochi \prg.
\end{itemize}
\end{chapgoals}

% ----------------------------------------------------------------------------
\section{Riscaldamento sulla macchina}

Una palla deve attraversare lentamente lo schermo. Un pixel per frame è troppo
veloce; zero pixel per frame non è un movimento. Alla palla serve una via di
mezzo: 0,375 pixel per frame. Ma il processore del kit è RV32IMAC: ha solo
istruzioni intere, e nessuna virgola mobile. Il Listato~\ref{lst:real:ball}
risolve il problema con un trucco.

\asmfile[caption={Tre ottavi di pixel per frame, con i soli interi.},label=lst:real:ball]{code/l06/s04.S}

La posizione è memorizzata \emph{moltiplicata per 256}, e viene divisa per 256
solo quando la palla viene disegnata. In unità di $1/256$ di pixel la velocità
è $0,375\times 256 = 96$, un onesto intero.

\begin{prgtry}
Eseguite il programma con \code{SPEED\_Q8} uguale a 96, poi 16, poi 1.
\end{prgtry}

Il trucco ha un nome, \term{virgola fissa}. Questo capitolo dice da dove viene
e quando non basta, nel qual caso subentra la risposta generale, la
\term{virgola mobile}. Prima dell'una e dell'altra servono altre due codifiche
degli interi con segno, perché la virgola mobile è costruita con esse.

% ----------------------------------------------------------------------------
\section{Altre codifiche per gli interi con segno}

\subsection{Complemento a uno}

\begin{definition}{Complemento a uno}
Su $n$ bit, un numero non negativo $0\le x<2^{n-1}$ è rappresentato da sé
stesso, $X=x$; un numero negativo è rappresentato da $X = 2^{n}-1-|x|$.
\end{definition}

L'intervallo va da $-(2^{n-1}-1)$ a $+(2^{n-1}-1)$. Lo zero ha due
rappresentazioni, \texttt{0000} e \texttt{1111}, quindi si rappresentano
$2^{n}-1$ numeri, come in segno e modulo. La ragione per usarlo è che il cambio
di segno è banale: poiché $2^{n}-1$ è una stringa di uni, sottrarre da esso
complementa ogni bit,
\[ -x \;=\; \mathtt{not}\ x. \]
Quindi $6=\texttt{0110}$ e $-6=\texttt{1001}$; $5=\texttt{0101}$ e
$-5=\texttt{1010}$; $1=\texttt{0001}$ e $-1=\texttt{1110}$. La \ins{not} bit a
bit di \riscv calcola esattamente questo.

L'aritmetica in complemento a uno lavora modulo $2^{n}-1$ invece che modulo
$2^{n}$, e questo costa una correzione.

\begin{center}\ttfamily
\begin{tabular}{@{}r@{\quad}r@{}}
\multicolumn{2}{c}{\textsf{\small Nessuna correzione}}\\
3 + & 0011 +\\ -4 = & 1011 =\\\hline -1 & 1110\\ &
\end{tabular}\qquad\qquad
\begin{tabular}{@{}r@{\quad}r@{}}
\multicolumn{2}{c}{\textsf{\small Due negativi}}\\
-2 + & 1101 +\\ -3 = & 1100 =\\\hline & \textcolor{prgred}{1}\,1001 + 1 =\\ -5 & 1010
\end{tabular}\qquad\qquad
\begin{tabular}{@{}r@{\quad}r@{}}
\multicolumn{2}{c}{\textsf{\small Risultato positivo}}\\
5 + & 0101 +\\ -2 = & 1101 =\\\hline & \textcolor{prgred}{1}\,0010 + 1 =\\ 3 & 0011
\end{tabular}
\end{center}

Ogni volta che c'è un riporto uscente dal bit più significativo, al risultato
va sommato 1: il \term{riporto circolare} (end-around carry). La correzione
serve quando entrambi gli operandi sono negativi, oppure quando gli operandi
hanno segni opposti e la somma è positiva. Un sommatore e un complementatore
bastano ancora; ma il passo in più e il doppio zero hanno fatto vincere il
complemento a due.

\subsection{Rappresentazione in eccesso-$k$}

\begin{definition}{Eccesso-$k$}
Nella rappresentazione in \term{eccesso-$k$} ogni numero è traslato verso
l'alto di una costante $k$, il \term{bias}: $X = x+k$. Tipicamente
$k=2^{n-1}$.
\end{definition}

Su quattro bit con $k=8$ la configurazione \texttt{0000} significa $-8$,
\texttt{1000} significa 0 e \texttt{1111} significa $+7$. La proprietà decisiva
è che l'ordine è conservato: se $x_1>x_2$ allora $X_1>X_2$, quindi i numeri con
segno si confrontano come se fossero naturali. Il prezzo si paga
nell'aritmetica: una somma va corretta sottraendo $k$, perché
$X_1+X_2 = (x_1+x_2)+2k$, e moltiplicazione e divisione sono molto più scomode
che con i complementi. L'eccesso-$k$ è quindi preferito dove servono solo somme
e confronti, ed è esattamente la situazione dell'esponente di un numero in
virgola mobile.

\subsection{Quattro bit, cinque letture}

La Tabella~\ref{tab:real:five} completa la tabella del
Capitolo~\ref{ch:integers}. Segno e modulo e complemento a uno hanno due zeri.
Complemento a due ed eccesso-8 coprono lo stesso intervallo, da $-8$ a $+7$,
con un solo zero; differiscono soltanto nel bit più significativo, che è
invertito. Solo l'eccesso-$k$ mantiene le configurazioni nello stesso ordine
dei numeri.

\begin{table}[htb]
\centering\small
\caption{Quattro bit, cinque letture. S-M: segno e modulo; C1, C2: complemento a uno e a due; Ecc-8: eccesso-8.}\label{tab:real:five}
\setlength{\tabcolsep}{4.5pt}
\begin{tabular}{@{}cccccc@{\qquad}cccccc@{}}
\toprule
\textbf{Bit} & \textbf{Nat.} & \textbf{S-M} & \textbf{C1} & \textbf{C2} & \textbf{Ecc-8} &
\textbf{Bit} & \textbf{Nat.} & \textbf{S-M} & \textbf{C1} & \textbf{C2} & \textbf{Ecc-8}\\
\midrule
\texttt{0000} & 0 & $+0$ & $+0$ & 0 & $-8$ & \texttt{1000} & 8 & $-0$ & $-7$ & $-8$ & 0\\
\texttt{0001} & 1 & $+1$ & $+1$ & $+1$ & $-7$ & \texttt{1001} & 9 & $-1$ & $-6$ & $-7$ & $+1$\\
\texttt{0010} & 2 & $+2$ & $+2$ & $+2$ & $-6$ & \texttt{1010} & 10 & $-2$ & $-5$ & $-6$ & $+2$\\
\texttt{0011} & 3 & $+3$ & $+3$ & $+3$ & $-5$ & \texttt{1011} & 11 & $-3$ & $-4$ & $-5$ & $+3$\\
\texttt{0100} & 4 & $+4$ & $+4$ & $+4$ & $-4$ & \texttt{1100} & 12 & $-4$ & $-3$ & $-4$ & $+4$\\
\texttt{0101} & 5 & $+5$ & $+5$ & $+5$ & $-3$ & \texttt{1101} & 13 & $-5$ & $-2$ & $-3$ & $+5$\\
\texttt{0110} & 6 & $+6$ & $+6$ & $+6$ & $-2$ & \texttt{1110} & 14 & $-6$ & $-1$ & $-2$ & $+6$\\
\texttt{0111} & 7 & $+7$ & $+7$ & $+7$ & $-1$ & \texttt{1111} & 15 & $-7$ & $-0$ & $-1$ & $+7$\\
\bottomrule
\end{tabular}
\end{table}

% ----------------------------------------------------------------------------
\section{Virgola fissa}

\subsection{Numeri in virgola fissa}

In virgola fissa la parte intera e la parte frazionaria di un numero sono
memorizzate una accanto all'altra in una sola parola. La posizione della
virgola è fissa e implicita: non è memorizzata, la conosce il programma.

\begin{definition}{Virgola fissa}
Con $f$ bit frazionari un numero $x$ è memorizzato come l'intero
$X = x\cdot 2^{f}$ (arrotondato). La \term{risoluzione} è $2^{-f}$. La
notazione Q$m.f$ indica un formato con $m$ bit interi e $f$ bit frazionari.
\end{definition}

Di solito l'intera parola è in complemento a due. La virgola fissa è il modo in
cui giochi, audio e software di controllo trattano le frazioni sui processori
privi di un'unità in virgola mobile. È anche il modo in cui andrebbe
memorizzato il denaro: come numero intero di centesimi.

\subsection{La palla, bit per bit: Q24.8}

La palla del riscaldamento è nel formato Q24.8
(Figura~\ref{fig:real:q248}). Quando la sua variabile contiene
$\hex{00002A60}=10848$, la palla si trova a $10848/256 = 42,375$ pixel: i 24
bit alti, \hex{00002A}, sono i 42 pixel interi e gli 8 bassi, \hex{60}, sono
$96/256 = 0,375$.

\begin{figure}[htb]
\centering
\begin{bytefield}[bitwidth=0.029\textwidth,bitheight=8mm]{32}
\bitheader[endianness=big]{0,7,8,31}\\
\bitbox{24}[bgcolor=prgaccent!20]{\texttt{0000 0000 0000 0000 0010 1010}} & \bitbox{8}[bgcolor=prgamber!30]{\texttt{0110 0000}}\\
\bitbox[]{24}{\footnotesize parte intera: 24 bit\\\texttt{0x00002A} = 42 pixel} & \bitbox[]{8}{\footnotesize frazione: 8 bit\\\texttt{0x60} = 96/256}
\end{bytefield}
\caption{\code{ball\_x\_q8 = 0x00002A60}: la palla si trova a 42,375 pixel.}
\label{fig:real:q248}
\end{figure}

Dopo un altro frame il valore è \hex{2AC0} (42,75), poi \hex{2B20} (43,125):
la palla attraversa il confine di un pixel ogni due o tre frame, e in media
avanza di 0,375.

L'aritmetica in virgola fissa è aritmetica intera con un po' di attenzione
(Tabella~\ref{tab:real:ops}). Somme e differenze non ne richiedono, perché le
virgole dei due operandi sono allineate. Il prodotto di due numeri scalati di
$2^{f}$ è scalato di $2^{2f}$, e va riportato indietro con uno scorrimento di
$f$. Lo scorrimento che riporta ai pixel interi è quello \emph{aritmetico},
\ins{srai}, che conserva il segno (Capitolo~\ref{ch:integers}).

\begin{table}[htb]
\centering\small
\caption{Le operazioni in virgola fissa.}\label{tab:real:ops}
\begin{tabular}{@{}lll@{}}
\toprule
\textbf{Operazione} & \textbf{In virgola fissa} & \textbf{\riscv}\\
\midrule
somma, sottrazione & come tra interi: le virgole sono allineate & \ins{add}, \ins{sub}, \ins{addi}\\
ai pixel interi & scorrimento a destra di $f$, conservando il segno & \code{srai a0, t1, 8}\\
prodotto & moltiplicare, poi scorrere a destra di $f$ & \ins{mul}, poi \ins{srai}\\
\bottomrule
\end{tabular}
\end{table}

\subsection{Numeri razionali: l'errore di approssimazione}

I numeri reali sono densi: anche in un intervallo limitato, non ogni numero si
può rappresentare con un numero finito di bit. A ogni numero si associa quindi
il più vicino tra quelli rappresentabili. In Q24.8 il numero 0,1 diventa
$26/256 = 0,1015625$, perché $0,1\times 256 = 25,6$ viene arrotondato a 26: un
errore di circa l'1,6\%.

L'errore di approssimazione è al massimo metà della risoluzione: $1/512$ con
otto bit frazionari. L'\term{underflow} è il caso di un numero vicino a zero,
ma non zero, che viene rappresentato da zero: con una velocità di 0,001 pixel
per frame la costante verrebbe arrotondata a 0, e la palla non si muoverebbe
mai.

\begin{prgconcept}
In virgola fissa l'errore \emph{assoluto} è lo stesso ovunque sulla retta dei
numeri. La risoluzione di otto bit frazionari è $1/256$: nulla di più piccolo
si distingue da zero.
\end{prgconcept}

% ----------------------------------------------------------------------------
\section{Virgola mobile}

La virgola fissa spende tutti i suoi bit su un unico intervallo fissato. La
virgola mobile ne spende alcuni per dire \emph{dove} si trova l'intervallo:
scambia precisione con ampiezza dell'intervallo.

\subsection{Numeri reali in virgola mobile}

\begin{definition}{Virgola mobile}
Un numero reale $x$ è rappresentato in virgola mobile da una coppia $(m,e)$
tale che
\[ x = m\cdot b^{\,e}, \]
dove $m$ è la \term{mantissa}, $e$ l'\term{esponente} e $b$ la base.
\end{definition}

La virgola è «mobile» perché è l'esponente a dire dove si trova. È la notazione
scientifica del laboratorio di fisica. Infinite coppie rappresentano lo stesso
numero,
\[
346,09801 \;=\; (346,09801;\,0) \;=\; (346098,01;\,-3) \;=\; (0,034609801;\,4),
\]
e la \term{normalizzazione} sceglie un unico rappresentante: la virgola va
subito dopo la prima cifra significativa, $(3,4609801;\,2)$.

\subsection{L'intervallo di un piccolo formato decimale}

Si consideri un formato in base 10 con quattro cifre per la mantissa e due per
l'esponente (Figura~\ref{fig:real:line}). Rappresenta i numeri da
$-9,999\times 10^{99}$ a $-1,000\times 10^{-99}$, lo zero, e i numeri da
$+1,000\times 10^{-99}$ a $+9,999\times 10^{99}$. Si ha un \term{overflow}
quando il modulo è troppo grande per l'esponente; un \term{underflow} quando è
troppo piccolo, e il numero collassa a zero.

\begin{figure}[htb]
\centering
\begin{tikzpicture}[x=1cm,y=1cm]
  \draw[<->,thick,prgink] (-7,0) -- (7,0);
  \fill[prgaccent!30] (-5.6,-0.22) rectangle (-1.0,0.22);
  \fill[prgaccent!30] (1.0,-0.22) rectangle (5.6,0.22);
  \fill[prgred!25] (-7,-0.22) rectangle (-5.6,0.22);
  \fill[prgred!25] (5.6,-0.22) rectangle (7,0.22);
  \fill[prgamber!35] (-1.0,-0.22) rectangle (1.0,0.22);
  \draw[thick] (0,-0.3) -- (0,0.3) node[above,lbl]{0};
  \foreach \x/\l in {-5.6/{$-9,999\times10^{99}$},-1.0/{$-1,000\times10^{-99}$},1.0/{$1,000\times10^{-99}$},5.6/{$9,999\times10^{99}$}}
    \draw (\x,-0.3) -- (\x,0.3) node[above,font=\scriptsize]{\l};
  \node[lbl,text=prgred] at (-6.3,-0.6) {overflow};
  \node[lbl,text=prgred] at (6.3,-0.6) {overflow};
  \node[lbl,text=prgamber] at (0,-0.6) {underflow};
  \node[lbl,text=prgaccent] at (-3.3,-0.6) {rappresentabili};
  \node[lbl,text=prgaccent] at (3.3,-0.6) {rappresentabili};
\end{tikzpicture}
\caption{Un formato decimale con quattro cifre di mantissa e due di esponente
(non in scala).}
\label{fig:real:line}
\end{figure}

\subsection{Approssimazione: errore assoluto ed errore relativo}

In un formato del genere l'errore assoluto non è costante. Vicino a zero la
distanza tra due numeri consecutivi è
$1,001\times 10^{-99}-1,000\times 10^{-99} = 0,001\times 10^{-99}$; vicino al
limite superiore è
$9,999\times 10^{99}-9,998\times 10^{99} = 0,001\times 10^{99}$. L'errore
assoluto è piccolo sui numeri piccoli e grande sui numeri grandi. Ciò che è
costante lungo tutto l'asse è l'errore \emph{relativo}, perché il numero di
cifre significative, quattro, è lo stesso ovunque.

\begin{prgconcept}
Virgola fissa: errore assoluto costante, errore relativo che esplode vicino a
zero. Virgola mobile: errore relativo costante, errore assoluto che cresce con
il numero.
\end{prgconcept}

\subsection{IEEE 754 a precisione singola}

Fino al 1985 ogni costruttore aveva un proprio formato in virgola mobile, e i
programmi davano risultati diversi su macchine diverse. Lo standard IEEE~754
(1985, rivisto nel 2008 e nel 2019) vi ha posto fine \citep{ieee754}. Il suo
formato a precisione singola divide 32 bit in tre campi
(Figura~\ref{fig:real:ieee}).

\begin{figure}[htb]
\centering
\begin{bytefield}[bitwidth=0.029\textwidth,bitheight=8mm]{32}
\bitheader[endianness=big]{0,22,23,30,31}\\
\bitbox{1}[bgcolor=prgred!30]{S} & \bitbox{8}[bgcolor=prgamber!30]{E} & \bitbox{23}[bgcolor=prgaccent!20]{F}\\
\bitbox[]{1}{\scriptsize segno} & \bitbox[]{8}{\footnotesize esponente: 8 bit\\eccesso-127} & \bitbox[]{23}{\footnotesize frazione: 23 bit\\modulo, dopo l'1 nascosto}
\end{bytefield}
\[ \text{valore} \;=\; (-1)^{S}\times 1,F\times 2^{\,E-127} \]
\caption{IEEE 754 a precisione singola.}
\label{fig:real:ieee}
\end{figure}

\begin{itemize}
\item La mantissa è in segno e modulo: un bit di segno e 23 bit di modulo.
\item Una mantissa binaria normalizzata comincia sempre con 1. Quel bit è
  \term{nascosto}: non viene memorizzato, e si guadagna un bit di precisione.
\item L'esponente è in eccesso-127 (non 128): gli esponenti veri vanno da
  $-126$ a $+127$. Grazie alla codifica in eccesso, e alla posizione
  dell'esponente prima della frazione, due numeri positivi in virgola mobile si
  possono confrontare come se fossero naturali.
\item Le configurazioni $E=0$ ed $E=255$ sono riservate: zero e numeri
  denormalizzati, infinito e NaN.
\end{itemize}

\begin{example}{$-5,125$ in IEEE 754}
\begin{enumerate}[nosep]
\item Si convertono le due parti: $5=101_2$ e $0,125=1/8=0,001_2$, quindi
  $5,125=101,001_2$.
\item Si normalizza: $101,001_2 = 1,01001_2\times 2^{2}$.
\item Segno $S=1$; esponente $E = 2+127 = 129 = \texttt{1000\,0001}$.
\item Frazione $F=\texttt{01001}$ seguita da zeri: l'1 iniziale è nascosto.
\end{enumerate}
\[
\underbrace{\texttt{1}}_{S}\;\underbrace{\texttt{1000\,0001}}_{E}\;
\underbrace{\texttt{0100\,1000\,0000\,0000\,0000\,000}}_{F}
\]
Raggruppato a nibble, \texttt{1100 0000 1010 0100 0000 \ldots} è
\hex{C0A40000}.
\end{example}

\subsection{IEEE 754 in sintesi}

La Tabella~\ref{tab:real:formats} confronta i due formati di uso comune, e la
Tabella~\ref{tab:real:special} elenca il significato di tutte le
configurazioni.

\begin{table}[htb]
\centering\small
\caption{IEEE 754 a precisione singola e doppia.}\label{tab:real:formats}
\begin{tabular}{@{}lll@{}}
\toprule
 & \textbf{Precisione singola} & \textbf{Precisione doppia}\\
\midrule
Bit totali & 32 & 64\\
Bit di segno & 1 & 1\\
Bit di esponente & 8 & 11\\
Bit di frazione & 23 & 52\\
Codifica dell'esponente & eccesso-127 & eccesso-1023\\
Intervallo dell'esponente & $-126\ldots+127$ & $-1022\ldots+1023$\\
Minimo normalizzato & $2^{-126}\approx 1,2\times 10^{-38}$ & $2^{-1022}\approx 2,2\times 10^{-308}$\\
Massimo & circa $2^{128}\approx 3,4\times 10^{38}$ & circa $2^{1024}\approx 1,8\times 10^{308}$\\
Minimo denormalizzato & circa $1,4\times 10^{-45}$ & circa $4,9\times 10^{-324}$\\
Cifre decimali & circa 7 & circa 16\\
\bottomrule
\end{tabular}
\end{table}

\begin{table}[htb]
\centering\small
\caption{IEEE 754: che cosa significano le configurazioni.}\label{tab:real:special}
\begin{tabular}{@{}lllll@{}}
\toprule
\textbf{Tipo} & \textbf{Segno} & \textbf{Esponente $E$} & \textbf{Frazione $F$} & \textbf{Valore}\\
\midrule
Normalizzato & 0 o 1 & $0<E<\max$ & qualunque & $(-1)^{S}\times 1,F\times 2^{E-\mathit{bias}}$\\
Denormalizzato & 0 o 1 & tutti zeri & non zero & $(-1)^{S}\times 0,F\times 2^{1-\mathit{bias}}$\\
Zero & 0 o 1 & tutti zeri & tutti zeri & $+0$ o $-0$\\
Infinito & 0 o 1 & tutti uni & tutti zeri & $+\infty$ o $-\infty$\\
Not a Number & 0 o 1 & tutti uni & non zero & NaN: un risultato non valido\\
\bottomrule
\end{tabular}
\end{table}

I numeri \term{denormalizzati} (o subnormali) riempiono il vuoto
dell'underflow attorno allo zero con una precisione che diminuisce
gradualmente. La divisione $1/0$ dà infinito; $0/0$ e la radice quadrata di
$-1$ danno NaN. NaN non è uguale a nulla, nemmeno a sé stesso.

\subsection{La virgola mobile non è l'insieme dei numeri reali}

Un numero in virgola mobile è un numero razionale scelto da un insieme finito,
e l'aritmetica su tali numeri obbedisce a leggi che non sono quelle del campo
reale \citep{goldberg1991}.

\begin{itemize}
\item 0,1 non ha un'espansione binaria finita: è memorizzato in modo
  approssimato, come 0,8 nel Capitolo~\ref{ch:integers}.
\item Due risultati in virgola mobile non vanno mai confrontati per
  uguaglianza: se ne confronta la differenza con una tolleranza
  (Listato~\ref{lst:real:c}).
\item L'addizione non è associativa: $(a+b)+c$ può differire da $a+(b+c)$.
\item Sommare un numero molto piccolo a uno molto grande può lasciare invariato
  quello grande.
\item Per il denaro e per le posizioni di un gioco, gli interi scalati sono
  spesso lo strumento giusto.
\end{itemize}

\cfile[caption={La somma più famosa dell'informatica.},label=lst:real:c,numbers=none]{code/l06/s22.c}

\subsection{Oltre la singola e la doppia}

La scelta di una rappresentazione è una decisione ingegneristica che si prende
ancora oggi. La \emph{mezza precisione}, aggiunta allo standard nel 2008, ha
sedici bit: 1 di segno, 5 di esponente, 10 di frazione; dà circa tre cifre
decimali in metà della memoria della precisione singola. Conta perché le reti
neurali contengono miliardi di numeri: dimezzare ciascuno dimezza memoria,
banda ed energia. Molti modelli vanno oltre e vengono memorizzati, per
\term{quantizzazione}, come interi a 8 bit con un fattore di scala: di nuovo
virgola fissa, su scala industriale. I bit sono un bilancio; intervallo,
precisione e costo tirano in direzioni diverse, su una GPU come su una palla in
Q24.8.

% ----------------------------------------------------------------------------
\section{Ritorno alla macchina}

\subsection{La virgola mobile su RISC-V, e sul kit}

\riscv è modulare: una piccola base obbligatoria più estensioni facoltative
(Tabella~\ref{tab:real:rv}). Un processore \riscv con le estensioni F e D ha
altri trentadue registri, da \reg{f0} a \reg{f31}, per i valori in virgola
mobile, e istruzioni come \ins{fadd.s}, \ins{fmul.s}, \ins{flw} e \ins{fsw} che
lavorano su numeri IEEE~754; \ins{fcvt.w.s} converte in intero e
\ins{fcvt.s.w} da intero. Processori così si trovano in portatili, server e
supercalcolatori.

L'ESP32-C6 è RV32IMAC: non ha né F né D, e quindi nessuna istruzione in virgola
mobile. Un programma C può comunque usare il tipo \code{float}: il compilatore
chiama routine software come \code{\_\_addsf3}, che realizzano esattamente
l'algoritmo di questo capitolo con scorrimenti, maschere e somme intere, al
costo di decine o centinaia di istruzioni intere per operazione. I giochi \prg
usano interi e virgola fissa: veloci, abbastanza esatti, e visibili nel
debugger.

\begin{table}[htb]
\centering\small
\caption{Leggere il nome RV32IMAC.}\label{tab:real:rv}
\begin{tabularx}{\textwidth}{@{}llLc@{}}
\toprule
\textbf{Lettera} & \textbf{Estensione} & \textbf{Che cosa aggiunge} & \textbf{ESP32-C6}\\
\midrule
RV32I & insieme intero di base & registri a 32 bit, somma, logica, load e store, salti & sì\\
M & moltiplicazione e divisione & \ins{mul}, \ins{mulh}, \ins{div}, \ins{rem} e le loro forme senza segno & sì\\
A & atomiche & istruzioni per la condivisione sicura tra task & sì\\
C & compresse & codifiche a 16 bit delle istruzioni più comuni & sì\\
F & precisione singola & da \reg{f0} a \reg{f31} e istruzioni IEEE 754 a precisione singola & no\\
D & precisione doppia & istruzioni IEEE 754 a precisione doppia & no\\
\bottomrule
\end{tabularx}
\end{table}

Un microcontrollore sceglie le estensioni che gli servono. È una delle ragioni
per cui la stessa architettura scala da un chip che costa due euro a un
supercalcolatore.

% ----------------------------------------------------------------------------
\section{Verificate di aver capito}

\begin{quiz}{Quale codifica con segno mantiene le configurazioni di bit nello stesso ordine dei numeri?}
\item Segno e modulo
\item Complemento a uno
\item Complemento a due
\item Eccesso-$k$
\end{quiz}

\begin{quiz}{In virgola fissa Q24.8, quale numero reale rappresenta la parola \hex{00000180}?}
\item 384
\item 1,5
\item 1,8
\item 0,375
\end{quiz}

\begin{quiz}{Quale numero reale è la configurazione IEEE 754 a precisione singola \hex{40400000}?}
\item 2,0
\item 3,0
\item 4,0
\item 0,75
\end{quiz}

\begin{quiz}{In IEEE 754, un campo esponente di tutti uni con una frazione di tutti zeri significa:}
\item zero
\item un numero denormalizzato
\item infinito
\item NaN
\end{quiz}

\section{Esercizi}

\begin{exercise}[prima su carta, poi su \prg]
\begin{enumerate}[nosep]
\item Scrivete $-6$ su quattro bit in segno e modulo, complemento a uno,
  complemento a due ed eccesso-8.
\item Codificate 10,5 e 0,15625 in IEEE 754 a precisione singola; date il
  risultato in esadecimale.
\item Decodificate \hex{BF800000}.
\item Eseguite la palla con \code{SPEED\_Q8} uguale a 96, 16 e 1. Dopo quanti
  frame si sposta di un pixel in ciascun caso?
\item Date alla palla anche una velocità verticale, e fatela rimbalzare:
  cambiate di segno la velocità con \ins{neg} quando raggiunge un bordo.
\end{enumerate}
\expected{Su carta: \texttt{1110}, \texttt{1001}, \texttt{1010},
\texttt{0010}; \hex{41280000} e \hex{3E200000}; $-1,0$. Su \prg: la palla
avanza di un pixel ogni 2,67, 16 e 256 frame; con il rimbalzo si muove in
diagonale e resta sullo schermo.}
\end{exercise}

\section{Glossario}

\begin{glossarytable}
one's complement & complemento a uno (alla base diminuita)\\
excess-$k$ (biased) & eccesso-$k$\\
fixed point & virgola fissa\\
floating point & virgola mobile\\
single / double precision & singola / doppia precisione\\
rounding & arrotondamento\\
mantissa / fraction & mantissa\\
exponent & esponente\\
normalisation & normalizzazione\\
hidden bit & bit nascosto\\
denormalised & denormalizzato\\
Not a Number (NaN) & non un numero (NaN)\\
\end{glossarytable}

\section{Conclusioni e letture}

Gli interi hanno più rappresentazioni: 8, 16, 32 o 64 bit, con segno o senza,
complementi o eccesso. La virgola fissa memorizza una frazione come intero
scalato: errore assoluto costante, la velocità degli interi. La virgola mobile
memorizza una mantissa e un esponente: un intervallo enorme, errore relativo
costante. IEEE~754 è lo standard universale per la virgola mobile, con
configurazioni speciali per zero, infinito e NaN. Su un processore senza
l'estensione F, la virgola fissa è la scelta naturale.

\paragraph{Letture.} \citet{patterson2020}, paragrafo~3.5;
\citet{goldberg1991}; \emph{Playing with RISC-V}, l'appendice sull'interfaccia
del runtime, per playfield e parallasse in Q8 \citep{prg32textbook}.

Il primo blocco è completo. Il prossimo capitolo mette al lavoro le sue quattro
lezioni in un programma vero: la prima parte di \vermone.